% Updated research design aligned with the completed 2026 World Cup dataset.
% This file is intended to be included with % Updated research design aligned with the completed 2026 World Cup dataset.
% This file is intended to be included with % Updated research design aligned with the completed 2026 World Cup dataset.
% This file is intended to be included with % Updated research design aligned with the completed 2026 World Cup dataset.
% This file is intended to be included with \input{sections/research_design}.
% Verified references use natbib; remaining legacy citations await metadata verification.

\section{Research Design}
\label{sec:research-design}

\subsection{Research objective and setting}

This study examines whether unusually large individual trades are associated with short-horizon movements in Polymarket-implied probabilities, whether additional flow-aligned movement follows between minutes 5 and 15 or instead reverses, and whether prices improve alignment with eventual market outcomes. The design is observational: it estimates conditional associations and does not claim that a trade caused a price change, conveyed private information, or manipulated a market. The primary outcome is described as a recorded-price adjustment, and the estimated coefficients as directional price associations.

The observation window spans 2026-06-01T00:00:00Z to 2026-08-01T00:00:00Z, with the end point excluded. The final universe consists of 360 binary markets: 312 home-win, away-win, or draw markets covering 104 FIFA World Cup matches, and 48 outright country-winner markets. The canonical trade sample contains 6,725,732 records from the public Polymarket Data API and 335,414 normalised proxy-wallet addresses. The price dataset contains 17,980,605 raw CLOB price-history points for 720 YES and NO tokens, requested at one-minute fidelity without interpolation or resampling. Match timing and outcomes come from the reviewed 104-event FIFA table. Independently obtained Dune Analytics market-level counts were used only to validate extraction completeness; Dune observations were not used to add, remove, or replace canonical trade observations.

\subsection{Research questions and hypotheses}

The study asks four questions. First, as a baseline check, how is signed trade flow associated with subsequent changes in a market's YES-equivalent probability? Second, and central to the dissertation, does unusually large net flow have a stronger association than ordinary net flow? Update incidence and conditional magnitude are examined as mechanism diagnostics rather than embedded in H2. Third, is large-trade flow associated with additional flow-aligned movement between minutes 5 and 15 or reversal over that interval, and do prices improve relative to the realised binary outcome? Fourth, does the association vary with time to kickoff and match phase?

H1 and H2 are primary. Within that pair, H1 is a baseline validation of the signed-flow construction, while H2, together with its extensive-versus-intensive-margin decomposition, is the dissertation's main empirical contribution. H3a, H3b and RQ4 are secondary hypotheses and a research question that interpret and add heterogeneity context to the H2 result rather than stand as co-equal contributions. The corresponding hypotheses and research question are:

\begin{description}
  \item[H1 (baseline).] Greater signed total trade flow is associated with a short-horizon probability change in the same direction. This result validates the signed-flow construction; it is not, by itself, the dissertation's central finding.
  \item[H2 (primary contribution).] Conditional on the common controls and fixed effects, signed flow from individual trades classified by a past-only, market-specific P99 threshold has a larger five-minute directional-price association than signed flow from ordinary trades. Whether the difference operates through update incidence or conditional update magnitude is evaluated as a mechanism decomposition rather than built into the hypothesis.
  \item[H3a.] Large-trade flow is associated with subsequent same-direction continuation from minute five to minute fifteen rather than reversal over that later interval.
  \item[H3b.] Large-trade-associated price movements improve forecast accuracy relative to the eventual resolved outcome, measured by Brier-score improvement, relative to ordinary-trade-associated movements.
\end{description}

The literature reviewed in Section~\ref{sec:literature-review} does not yield a strong ex-ante directional prediction for the shape of pre-kickoff timing heterogeneity, so this dimension is treated as a research question rather than forced into a falsifiable directional hypothesis:

\begin{description}
  \item[RQ4.] How does the large-versus-ordinary flow--price association vary across pre-specified time-to-kickoff bands, and between the pre-kickoff and post-kickoff/pre-resolution phases?
\end{description}

This formulation follows the microstructure literature on informed trading, trade size, and price discovery \citep{kyle1985continuous,easleyohara1987tradesize,hasbrouck1991information,madhavan1997prices}, while prediction-market research motivates the interpretation of prices as market-implied probabilities subject to institutional and belief-based qualifications \citep{wolferszitzewitz2004prediction,manski2006interpreting}. Manipulation and corrective-trading research motivates the distinction between temporary pressure and information \citep{allengale1992manipulation,camerer1998racetrack,hanson2006aggregation}.

\subsection{Observation units and sample separation}

The primary observation unit is the market-minute. For a minute beginning at
$t$, flow contains trades timestamped in $[t,t+1)$. The baseline is the latest
observed YES-equivalent price strictly before $t$, and each post-flow outcome is
the first observed price at or after the designated horizon. This ordering
prevents the baseline from incorporating a price observation recorded after a
trade in the focal minute.

The match and outright families are estimated separately. The primary match specification uses pre-kickoff observations. Post-kickoff/pre-resolution observations are analysed separately because match events generate fundamental information; this phase is not labelled purely in-play because actual final-whistle timestamps are unavailable. Outright markets form a separate panel because their horizon and information environment differ from single-match contracts. Trades at or after recorded resolution remain in the method-neutral audit panel but are excluded from model samples.

\subsection{Direction and large-trade definition}

Trade direction is expressed in YES-probability terms. BUY YES and SELL NO receive a sign of $+1$; SELL YES and BUY NO receive a sign of $-1$. For trade $i$,

\begin{equation}
  \text{SignedValue}_{i}=\text{Direction}_{i}\times
  \left(\text{ExecutedTokenPrice}_{i}\times\text{SizeShares}_{i}\right).
\end{equation}

A YES-token history is retained as observed, whereas a NO-token price is converted to one minus the NO price. YES and NO histories are not treated as independent markets. The YES history is primary and the NO-equivalent series is used for complementarity and data-quality checks.

The CLOB history documentation labels each observation only as a timestamp and
price value; it does not state that the value is a contemporaneous midpoint or
pre-trade quote. The primary dependent variable is therefore termed a
recorded-price adjustment. A robustness series instead uses the last executed
YES-equivalent trade price strictly before the focal minute and the first
executed trade price at or after the target horizon, with 60-, 300-, and
900-second matching tolerances. The focal minute is excluded from the
executed-price baseline.

For each binary market, trade value is executed token price multiplied by
shares. The primary expanding-history threshold for UTC day $d$ is the linearly
interpolated P99 of unsigned USDC trade value in the same market from 1 June to
the end of day $d-1$. A threshold becomes eligible after at least 1,000 prior
trades. The principal definition robustness instead uses the preceding seven
complete UTC days. A trade is large when its value is at or above the applicable
threshold; ties are included. Both definitions are past-only, so a trade is
never classified using later observations. Models are estimated separately and
on their common support. Past-only P95, P97.5, and P99.5 thresholds test whether
the finding depends on the P99 boundary. Signed-log flow is also replaced by
\texttt{asinh} flow and winsorised raw flow scaled in \$1,000 units;
standardised and common-value effects permit comparison despite the different
sparsity and scale of the large and ordinary components. Classification occurs before applying YES-equivalent
direction and does not represent wallet wealth, ownership, institutional status,
profitability, information, or intent.

\subsection{Price response, subsequent continuation, and forecast improvement}

Let $p^-_{m,t}$ be the latest available YES-equivalent CLOB price strictly
before $t$, and let $p^+_{m,t+h}$ be the first observed price at or after the
target $t+h$. Prices are not interpolated. For horizon $h$,

\begin{equation}
  \Delta p_{m,t}(h)=p^+_{m,t+h}-p^-_{m,t}.
\end{equation}

The primary pre-kickoff match horizon is five minutes; one, 15, and 30 minutes
trace the response path. Strict-0 alignment requires the baseline to occur less
than 60 seconds before $t$ and the outcome less than 60 seconds after its target;
strict-1 permits 120 seconds. No interpolation, resampling, or gap filling is
used. No pre-kickoff model window may cross actual kickoff.

Pre-trade placebo outcomes cover $[-15,0)$, $[-5,0)$, and $[-1,0)$ relative
to the focal flow minute. Post-trade lead--lag outcomes cover $+1$, $+5$,
$+15$, and $+30$ minutes. The seven focal tests are adjusted by Holm's method
within each threshold definition. A significant negative-horizon coefficient
would indicate anticipation, price-chasing, shared information, or residual
timing confounding rather than an effect of the later flow.

For $h_2>h_1$, subsequent movement is

\begin{equation}
  \text{Reversal}_{m,t}(h_1,h_2)=
  \operatorname{sign}(\text{NetSignedFlow}_{m,t})
  \left(p_{m,t+h_2}-p_{m,t+h_1}\right).
\end{equation}

A negative value denotes movement against the initial flow direction. A positive value is termed subsequent same-direction continuation, not strict persistence of the initial price movement: it can reflect delayed adjustment when no initial update occurred. A separate descriptive specification therefore conditions on a non-zero initial movement from minute zero to minute five. This conditional sample is outcome-selected and is not used as the primary H3a test. Reversal is not, by itself, proof of manipulation. Let $y_m\in\{0,1\}$ be the resolved outcome. Forecast improvement is measured by the change in Brier loss,

\begin{equation}
  \text{BrierImprovement}_{m,t}(h)=
  (y_m-p_{m,t})^2-(y_m-p_{m,t+h})^2.
\end{equation}

A positive value means that the later probability is closer to the realised outcome. This avoids an inappropriate unconditional 50 per cent accuracy benchmark. Movement to final zero-or-one settlement is not labelled permanent impact because settlement is mechanically determined by resolution.

\subsection{Empirical specifications}

The primary specification regresses the strictly aligned five-minute
pre-kickoff price change on signed-log flow. H1 uses total net signed flow. H2
decomposes flow into past-only P99 large-trade and ordinary-trade components and
directly tests $H_0:\beta_{large}=\beta_{ordinary}$. Primary controls are the
pre-period probability, lagged 30-minute price change, log lagged 30-minute
volume, and log minutes to kickoff. Models include market and unique UTC
calendar-hour fixed effects. Standard errors are clustered by the 104 match
events.

Two additional specifications address timing and shared match information.
First, the pre-trade placebo and post-trade lead--lag models estimate the same
coefficient contrast at all seven frozen horizons. Second, market fixed effects
are combined with event-by-minute and event-by-five-minute fixed effects, so
identification comes from related home-win, draw, and away-win contracts for the
same match in the same time cell. Cells containing only one observed market do
not identify this comparison and are excluded.

These event-by-time fixed effects absorb information common to related
contracts within the same match-time cell, thereby reducing common match-time
confounding. They cannot eliminate contract-specific responses to
contemporaneous match news. The resulting estimates therefore remain
associational.

The extensive margin is estimated first as a linear-probability model for any
non-zero five-minute update and then as a market-fixed-effects logit robustness
model. Markets without within-market outcome variation carry no finite logit
slope information and are excluded from that robustness sample. The intensive
margin conditions on a non-zero update and is treated as descriptive because it
selects on the realised outcome.

H3a is estimated on movement from minute 5 to minute 15; H3b is estimated on the associated Brier-score improvement. RQ4 uses five mutually exclusive pre-kickoff timing bands and a pooled five-minute pre-/post-kickoff phase-interaction model. The phase comparison uses controls common to both phases and tests the change in the large-minus-ordinary coefficient gap. Because RQ4 is exploratory, no monotonic pattern toward kickoff is imposed or expected a priori.

Primary standard errors are clustered by match event because the home-win, away-win, and draw contracts for the same FIFA match are related. Robustness considers market clustering, two-way event and calendar-hour clustering where supported, and non-overlapping five-minute observations. The clustering rule and number of clusters are stated for every table.

\subsection{Secondary wallet-heterogeneity design}
\label{sec:wallet-heterogeneity-design}

Secondary wallet tests distinguish an unusually large individual trade from
past observed-wallet activity and cumulative gross trading volume. High activity
is classified daily using a rolling seven-day top one per cent with minimum
history requirements; expanding history is a robustness definition. Four
large/ordinary by high/other activity cells are compared using price and
sequence outcomes. Complete definitions, HA1--HA4, tests, and multiplicity rules
are reported in Appendix~\ref{sec:wallet-methods-appendix}.

\subsection{Past-only whale-wallet design}
\label{sec:whale-wallet-design}

The past-only whale proxy instead ranks wallets by cumulative gross trading volume before
each evaluation day and imposes no minimum trade-frequency requirement. Rolling
windows, alternative percentiles, and continuous standardised cumulative volume
provide robustness. These measures describe prior observed gross trading volume,
not ownership, wealth, institutional status, or intent. Detailed WH1--WH4 and
outcome definitions are also provided in Appendix~\ref{sec:wallet-methods-appendix}.

\subsection{Identification, robustness, and data limitations}

Large trades may respond to public news or match events, creating simultaneity
and selection. Fixed effects, strict temporal ordering, pre-trade placebos, and
event-by-time comparisons reduce specific forms of confounding but do not
establish causality. The empirical language therefore remains associative.

\subsection{Reproducibility}

The public repository documents the final analysis workflow and contains a
README, stage-specific dependency files, the data dictionary and collection
report, analysis scripts, the principal non-confidential v5 and v6 output
files, SHA-256 manifests for 360 trade partitions and 360 price partitions,
and the complete LaTeX source. Large source and regression-ready datasets are
not redistributed in the repository. Their hash manifests identify the local
files used for the analysis after the public-source data have been collected.

The repository version corresponding to this submitted dissertation is
available at:
\url{https://github.com/andrew1004aa/polymarket-world-cup-2026/tree/v1.0.1-submission}.
Files retained for transparency or development history are not, merely by
their presence, part of the final research design described in this
dissertation.

Detailed specifications for the supplementary concentration,
reversal-prediction, wallet-sequence, and anomaly-screening analyses are reported in
Appendix~\ref{sec:supplementary-integrity-methods}. Keeping those specifications
outside the main chapter preserves the distinction between the primary H1--H3
design and exploratory behavioural or integrity analyses.
.
% Verified references use natbib; remaining legacy citations await metadata verification.

\section{Research Design}
\label{sec:research-design}

\subsection{Research objective and setting}

This study examines whether unusually large individual trades are associated with short-horizon movements in Polymarket-implied probabilities, whether additional flow-aligned movement follows between minutes 5 and 15 or instead reverses, and whether prices improve alignment with eventual market outcomes. The design is observational: it estimates conditional associations and does not claim that a trade caused a price change, conveyed private information, or manipulated a market. The primary outcome is described as a recorded-price adjustment, and the estimated coefficients as directional price associations.

The observation window spans 2026-06-01T00:00:00Z to 2026-08-01T00:00:00Z, with the end point excluded. The final universe consists of 360 binary markets: 312 home-win, away-win, or draw markets covering 104 FIFA World Cup matches, and 48 outright country-winner markets. The canonical trade sample contains 6,725,732 records from the public Polymarket Data API and 335,414 normalised proxy-wallet addresses. The price dataset contains 17,980,605 raw CLOB price-history points for 720 YES and NO tokens, requested at one-minute fidelity without interpolation or resampling. Match timing and outcomes come from the reviewed 104-event FIFA table. Independently obtained Dune Analytics market-level counts were used only to validate extraction completeness; Dune observations were not used to add, remove, or replace canonical trade observations.

\subsection{Research questions and hypotheses}

The study asks four questions. First, as a baseline check, how is signed trade flow associated with subsequent changes in a market's YES-equivalent probability? Second, and central to the dissertation, does unusually large net flow have a stronger association than ordinary net flow? Update incidence and conditional magnitude are examined as mechanism diagnostics rather than embedded in H2. Third, is large-trade flow associated with additional flow-aligned movement between minutes 5 and 15 or reversal over that interval, and do prices improve relative to the realised binary outcome? Fourth, does the association vary with time to kickoff and match phase?

H1 and H2 are primary. Within that pair, H1 is a baseline validation of the signed-flow construction, while H2, together with its extensive-versus-intensive-margin decomposition, is the dissertation's main empirical contribution. H3a, H3b and RQ4 are secondary hypotheses and a research question that interpret and add heterogeneity context to the H2 result rather than stand as co-equal contributions. The corresponding hypotheses and research question are:

\begin{description}
  \item[H1 (baseline).] Greater signed total trade flow is associated with a short-horizon probability change in the same direction. This result validates the signed-flow construction; it is not, by itself, the dissertation's central finding.
  \item[H2 (primary contribution).] Conditional on the common controls and fixed effects, signed flow from individual trades classified by a past-only, market-specific P99 threshold has a larger five-minute directional-price association than signed flow from ordinary trades. Whether the difference operates through update incidence or conditional update magnitude is evaluated as a mechanism decomposition rather than built into the hypothesis.
  \item[H3a.] Large-trade flow is associated with subsequent same-direction continuation from minute five to minute fifteen rather than reversal over that later interval.
  \item[H3b.] Large-trade-associated price movements improve forecast accuracy relative to the eventual resolved outcome, measured by Brier-score improvement, relative to ordinary-trade-associated movements.
\end{description}

The literature reviewed in Section~\ref{sec:literature-review} does not yield a strong ex-ante directional prediction for the shape of pre-kickoff timing heterogeneity, so this dimension is treated as a research question rather than forced into a falsifiable directional hypothesis:

\begin{description}
  \item[RQ4.] How does the large-versus-ordinary flow--price association vary across pre-specified time-to-kickoff bands, and between the pre-kickoff and post-kickoff/pre-resolution phases?
\end{description}

This formulation follows the microstructure literature on informed trading, trade size, and price discovery \citep{kyle1985continuous,easleyohara1987tradesize,hasbrouck1991information,madhavan1997prices}, while prediction-market research motivates the interpretation of prices as market-implied probabilities subject to institutional and belief-based qualifications \citep{wolferszitzewitz2004prediction,manski2006interpreting}. Manipulation and corrective-trading research motivates the distinction between temporary pressure and information \citep{allengale1992manipulation,camerer1998racetrack,hanson2006aggregation}.

\subsection{Observation units and sample separation}

The primary observation unit is the market-minute. For a minute beginning at
$t$, flow contains trades timestamped in $[t,t+1)$. The baseline is the latest
observed YES-equivalent price strictly before $t$, and each post-flow outcome is
the first observed price at or after the designated horizon. This ordering
prevents the baseline from incorporating a price observation recorded after a
trade in the focal minute.

The match and outright families are estimated separately. The primary match specification uses pre-kickoff observations. Post-kickoff/pre-resolution observations are analysed separately because match events generate fundamental information; this phase is not labelled purely in-play because actual final-whistle timestamps are unavailable. Outright markets form a separate panel because their horizon and information environment differ from single-match contracts. Trades at or after recorded resolution remain in the method-neutral audit panel but are excluded from model samples.

\subsection{Direction and large-trade definition}

Trade direction is expressed in YES-probability terms. BUY YES and SELL NO receive a sign of $+1$; SELL YES and BUY NO receive a sign of $-1$. For trade $i$,

\begin{equation}
  \text{SignedValue}_{i}=\text{Direction}_{i}\times
  \left(\text{ExecutedTokenPrice}_{i}\times\text{SizeShares}_{i}\right).
\end{equation}

A YES-token history is retained as observed, whereas a NO-token price is converted to one minus the NO price. YES and NO histories are not treated as independent markets. The YES history is primary and the NO-equivalent series is used for complementarity and data-quality checks.

The CLOB history documentation labels each observation only as a timestamp and
price value; it does not state that the value is a contemporaneous midpoint or
pre-trade quote. The primary dependent variable is therefore termed a
recorded-price adjustment. A robustness series instead uses the last executed
YES-equivalent trade price strictly before the focal minute and the first
executed trade price at or after the target horizon, with 60-, 300-, and
900-second matching tolerances. The focal minute is excluded from the
executed-price baseline.

For each binary market, trade value is executed token price multiplied by
shares. The primary expanding-history threshold for UTC day $d$ is the linearly
interpolated P99 of unsigned USDC trade value in the same market from 1 June to
the end of day $d-1$. A threshold becomes eligible after at least 1,000 prior
trades. The principal definition robustness instead uses the preceding seven
complete UTC days. A trade is large when its value is at or above the applicable
threshold; ties are included. Both definitions are past-only, so a trade is
never classified using later observations. Models are estimated separately and
on their common support. Past-only P95, P97.5, and P99.5 thresholds test whether
the finding depends on the P99 boundary. Signed-log flow is also replaced by
\texttt{asinh} flow and winsorised raw flow scaled in \$1,000 units;
standardised and common-value effects permit comparison despite the different
sparsity and scale of the large and ordinary components. Classification occurs before applying YES-equivalent
direction and does not represent wallet wealth, ownership, institutional status,
profitability, information, or intent.

\subsection{Price response, subsequent continuation, and forecast improvement}

Let $p^-_{m,t}$ be the latest available YES-equivalent CLOB price strictly
before $t$, and let $p^+_{m,t+h}$ be the first observed price at or after the
target $t+h$. Prices are not interpolated. For horizon $h$,

\begin{equation}
  \Delta p_{m,t}(h)=p^+_{m,t+h}-p^-_{m,t}.
\end{equation}

The primary pre-kickoff match horizon is five minutes; one, 15, and 30 minutes
trace the response path. Strict-0 alignment requires the baseline to occur less
than 60 seconds before $t$ and the outcome less than 60 seconds after its target;
strict-1 permits 120 seconds. No interpolation, resampling, or gap filling is
used. No pre-kickoff model window may cross actual kickoff.

Pre-trade placebo outcomes cover $[-15,0)$, $[-5,0)$, and $[-1,0)$ relative
to the focal flow minute. Post-trade lead--lag outcomes cover $+1$, $+5$,
$+15$, and $+30$ minutes. The seven focal tests are adjusted by Holm's method
within each threshold definition. A significant negative-horizon coefficient
would indicate anticipation, price-chasing, shared information, or residual
timing confounding rather than an effect of the later flow.

For $h_2>h_1$, subsequent movement is

\begin{equation}
  \text{Reversal}_{m,t}(h_1,h_2)=
  \operatorname{sign}(\text{NetSignedFlow}_{m,t})
  \left(p_{m,t+h_2}-p_{m,t+h_1}\right).
\end{equation}

A negative value denotes movement against the initial flow direction. A positive value is termed subsequent same-direction continuation, not strict persistence of the initial price movement: it can reflect delayed adjustment when no initial update occurred. A separate descriptive specification therefore conditions on a non-zero initial movement from minute zero to minute five. This conditional sample is outcome-selected and is not used as the primary H3a test. Reversal is not, by itself, proof of manipulation. Let $y_m\in\{0,1\}$ be the resolved outcome. Forecast improvement is measured by the change in Brier loss,

\begin{equation}
  \text{BrierImprovement}_{m,t}(h)=
  (y_m-p_{m,t})^2-(y_m-p_{m,t+h})^2.
\end{equation}

A positive value means that the later probability is closer to the realised outcome. This avoids an inappropriate unconditional 50 per cent accuracy benchmark. Movement to final zero-or-one settlement is not labelled permanent impact because settlement is mechanically determined by resolution.

\subsection{Empirical specifications}

The primary specification regresses the strictly aligned five-minute
pre-kickoff price change on signed-log flow. H1 uses total net signed flow. H2
decomposes flow into past-only P99 large-trade and ordinary-trade components and
directly tests $H_0:\beta_{large}=\beta_{ordinary}$. Primary controls are the
pre-period probability, lagged 30-minute price change, log lagged 30-minute
volume, and log minutes to kickoff. Models include market and unique UTC
calendar-hour fixed effects. Standard errors are clustered by the 104 match
events.

Two additional specifications address timing and shared match information.
First, the pre-trade placebo and post-trade lead--lag models estimate the same
coefficient contrast at all seven frozen horizons. Second, market fixed effects
are combined with event-by-minute and event-by-five-minute fixed effects, so
identification comes from related home-win, draw, and away-win contracts for the
same match in the same time cell. Cells containing only one observed market do
not identify this comparison and are excluded.

These event-by-time fixed effects absorb information common to related
contracts within the same match-time cell, thereby reducing common match-time
confounding. They cannot eliminate contract-specific responses to
contemporaneous match news. The resulting estimates therefore remain
associational.

The extensive margin is estimated first as a linear-probability model for any
non-zero five-minute update and then as a market-fixed-effects logit robustness
model. Markets without within-market outcome variation carry no finite logit
slope information and are excluded from that robustness sample. The intensive
margin conditions on a non-zero update and is treated as descriptive because it
selects on the realised outcome.

H3a is estimated on movement from minute 5 to minute 15; H3b is estimated on the associated Brier-score improvement. RQ4 uses five mutually exclusive pre-kickoff timing bands and a pooled five-minute pre-/post-kickoff phase-interaction model. The phase comparison uses controls common to both phases and tests the change in the large-minus-ordinary coefficient gap. Because RQ4 is exploratory, no monotonic pattern toward kickoff is imposed or expected a priori.

Primary standard errors are clustered by match event because the home-win, away-win, and draw contracts for the same FIFA match are related. Robustness considers market clustering, two-way event and calendar-hour clustering where supported, and non-overlapping five-minute observations. The clustering rule and number of clusters are stated for every table.

\subsection{Secondary wallet-heterogeneity design}
\label{sec:wallet-heterogeneity-design}

Secondary wallet tests distinguish an unusually large individual trade from
past observed-wallet activity and cumulative gross trading volume. High activity
is classified daily using a rolling seven-day top one per cent with minimum
history requirements; expanding history is a robustness definition. Four
large/ordinary by high/other activity cells are compared using price and
sequence outcomes. Complete definitions, HA1--HA4, tests, and multiplicity rules
are reported in Appendix~\ref{sec:wallet-methods-appendix}.

\subsection{Past-only whale-wallet design}
\label{sec:whale-wallet-design}

The past-only whale proxy instead ranks wallets by cumulative gross trading volume before
each evaluation day and imposes no minimum trade-frequency requirement. Rolling
windows, alternative percentiles, and continuous standardised cumulative volume
provide robustness. These measures describe prior observed gross trading volume,
not ownership, wealth, institutional status, or intent. Detailed WH1--WH4 and
outcome definitions are also provided in Appendix~\ref{sec:wallet-methods-appendix}.

\subsection{Identification, robustness, and data limitations}

Large trades may respond to public news or match events, creating simultaneity
and selection. Fixed effects, strict temporal ordering, pre-trade placebos, and
event-by-time comparisons reduce specific forms of confounding but do not
establish causality. The empirical language therefore remains associative.

\subsection{Reproducibility}

The public repository documents the final analysis workflow and contains a
README, stage-specific dependency files, the data dictionary and collection
report, analysis scripts, the principal non-confidential v5 and v6 output
files, SHA-256 manifests for 360 trade partitions and 360 price partitions,
and the complete LaTeX source. Large source and regression-ready datasets are
not redistributed in the repository. Their hash manifests identify the local
files used for the analysis after the public-source data have been collected.

The repository version corresponding to this submitted dissertation is
available at:
\url{https://github.com/andrew1004aa/polymarket-world-cup-2026/tree/v1.0.1-submission}.
Files retained for transparency or development history are not, merely by
their presence, part of the final research design described in this
dissertation.

Detailed specifications for the supplementary concentration,
reversal-prediction, wallet-sequence, and anomaly-screening analyses are reported in
Appendix~\ref{sec:supplementary-integrity-methods}. Keeping those specifications
outside the main chapter preserves the distinction between the primary H1--H3
design and exploratory behavioural or integrity analyses.
.
% Verified references use natbib; remaining legacy citations await metadata verification.

\section{Research Design}
\label{sec:research-design}

\subsection{Research objective and setting}

This study examines whether unusually large individual trades are associated with short-horizon movements in Polymarket-implied probabilities, whether additional flow-aligned movement follows between minutes 5 and 15 or instead reverses, and whether prices improve alignment with eventual market outcomes. The design is observational: it estimates conditional associations and does not claim that a trade caused a price change, conveyed private information, or manipulated a market. The primary outcome is described as a recorded-price adjustment, and the estimated coefficients as directional price associations.

The observation window spans 2026-06-01T00:00:00Z to 2026-08-01T00:00:00Z, with the end point excluded. The final universe consists of 360 binary markets: 312 home-win, away-win, or draw markets covering 104 FIFA World Cup matches, and 48 outright country-winner markets. The canonical trade sample contains 6,725,732 records from the public Polymarket Data API and 335,414 normalised proxy-wallet addresses. The price dataset contains 17,980,605 raw CLOB price-history points for 720 YES and NO tokens, requested at one-minute fidelity without interpolation or resampling. Match timing and outcomes come from the reviewed 104-event FIFA table. Independently obtained Dune Analytics market-level counts were used only to validate extraction completeness; Dune observations were not used to add, remove, or replace canonical trade observations.

\subsection{Research questions and hypotheses}

The study asks four questions. First, as a baseline check, how is signed trade flow associated with subsequent changes in a market's YES-equivalent probability? Second, and central to the dissertation, does unusually large net flow have a stronger association than ordinary net flow? Update incidence and conditional magnitude are examined as mechanism diagnostics rather than embedded in H2. Third, is large-trade flow associated with additional flow-aligned movement between minutes 5 and 15 or reversal over that interval, and do prices improve relative to the realised binary outcome? Fourth, does the association vary with time to kickoff and match phase?

H1 and H2 are primary. Within that pair, H1 is a baseline validation of the signed-flow construction, while H2, together with its extensive-versus-intensive-margin decomposition, is the dissertation's main empirical contribution. H3a, H3b and RQ4 are secondary hypotheses and a research question that interpret and add heterogeneity context to the H2 result rather than stand as co-equal contributions. The corresponding hypotheses and research question are:

\begin{description}
  \item[H1 (baseline).] Greater signed total trade flow is associated with a short-horizon probability change in the same direction. This result validates the signed-flow construction; it is not, by itself, the dissertation's central finding.
  \item[H2 (primary contribution).] Conditional on the common controls and fixed effects, signed flow from individual trades classified by a past-only, market-specific P99 threshold has a larger five-minute directional-price association than signed flow from ordinary trades. Whether the difference operates through update incidence or conditional update magnitude is evaluated as a mechanism decomposition rather than built into the hypothesis.
  \item[H3a.] Large-trade flow is associated with subsequent same-direction continuation from minute five to minute fifteen rather than reversal over that later interval.
  \item[H3b.] Large-trade-associated price movements improve forecast accuracy relative to the eventual resolved outcome, measured by Brier-score improvement, relative to ordinary-trade-associated movements.
\end{description}

The literature reviewed in Section~\ref{sec:literature-review} does not yield a strong ex-ante directional prediction for the shape of pre-kickoff timing heterogeneity, so this dimension is treated as a research question rather than forced into a falsifiable directional hypothesis:

\begin{description}
  \item[RQ4.] How does the large-versus-ordinary flow--price association vary across pre-specified time-to-kickoff bands, and between the pre-kickoff and post-kickoff/pre-resolution phases?
\end{description}

This formulation follows the microstructure literature on informed trading, trade size, and price discovery \citep{kyle1985continuous,easleyohara1987tradesize,hasbrouck1991information,madhavan1997prices}, while prediction-market research motivates the interpretation of prices as market-implied probabilities subject to institutional and belief-based qualifications \citep{wolferszitzewitz2004prediction,manski2006interpreting}. Manipulation and corrective-trading research motivates the distinction between temporary pressure and information \citep{allengale1992manipulation,camerer1998racetrack,hanson2006aggregation}.

\subsection{Observation units and sample separation}

The primary observation unit is the market-minute. For a minute beginning at
$t$, flow contains trades timestamped in $[t,t+1)$. The baseline is the latest
observed YES-equivalent price strictly before $t$, and each post-flow outcome is
the first observed price at or after the designated horizon. This ordering
prevents the baseline from incorporating a price observation recorded after a
trade in the focal minute.

The match and outright families are estimated separately. The primary match specification uses pre-kickoff observations. Post-kickoff/pre-resolution observations are analysed separately because match events generate fundamental information; this phase is not labelled purely in-play because actual final-whistle timestamps are unavailable. Outright markets form a separate panel because their horizon and information environment differ from single-match contracts. Trades at or after recorded resolution remain in the method-neutral audit panel but are excluded from model samples.

\subsection{Direction and large-trade definition}

Trade direction is expressed in YES-probability terms. BUY YES and SELL NO receive a sign of $+1$; SELL YES and BUY NO receive a sign of $-1$. For trade $i$,

\begin{equation}
  \text{SignedValue}_{i}=\text{Direction}_{i}\times
  \left(\text{ExecutedTokenPrice}_{i}\times\text{SizeShares}_{i}\right).
\end{equation}

A YES-token history is retained as observed, whereas a NO-token price is converted to one minus the NO price. YES and NO histories are not treated as independent markets. The YES history is primary and the NO-equivalent series is used for complementarity and data-quality checks.

The CLOB history documentation labels each observation only as a timestamp and
price value; it does not state that the value is a contemporaneous midpoint or
pre-trade quote. The primary dependent variable is therefore termed a
recorded-price adjustment. A robustness series instead uses the last executed
YES-equivalent trade price strictly before the focal minute and the first
executed trade price at or after the target horizon, with 60-, 300-, and
900-second matching tolerances. The focal minute is excluded from the
executed-price baseline.

For each binary market, trade value is executed token price multiplied by
shares. The primary expanding-history threshold for UTC day $d$ is the linearly
interpolated P99 of unsigned USDC trade value in the same market from 1 June to
the end of day $d-1$. A threshold becomes eligible after at least 1,000 prior
trades. The principal definition robustness instead uses the preceding seven
complete UTC days. A trade is large when its value is at or above the applicable
threshold; ties are included. Both definitions are past-only, so a trade is
never classified using later observations. Models are estimated separately and
on their common support. Past-only P95, P97.5, and P99.5 thresholds test whether
the finding depends on the P99 boundary. Signed-log flow is also replaced by
\texttt{asinh} flow and winsorised raw flow scaled in \$1,000 units;
standardised and common-value effects permit comparison despite the different
sparsity and scale of the large and ordinary components. Classification occurs before applying YES-equivalent
direction and does not represent wallet wealth, ownership, institutional status,
profitability, information, or intent.

\subsection{Price response, subsequent continuation, and forecast improvement}

Let $p^-_{m,t}$ be the latest available YES-equivalent CLOB price strictly
before $t$, and let $p^+_{m,t+h}$ be the first observed price at or after the
target $t+h$. Prices are not interpolated. For horizon $h$,

\begin{equation}
  \Delta p_{m,t}(h)=p^+_{m,t+h}-p^-_{m,t}.
\end{equation}

The primary pre-kickoff match horizon is five minutes; one, 15, and 30 minutes
trace the response path. Strict-0 alignment requires the baseline to occur less
than 60 seconds before $t$ and the outcome less than 60 seconds after its target;
strict-1 permits 120 seconds. No interpolation, resampling, or gap filling is
used. No pre-kickoff model window may cross actual kickoff.

Pre-trade placebo outcomes cover $[-15,0)$, $[-5,0)$, and $[-1,0)$ relative
to the focal flow minute. Post-trade lead--lag outcomes cover $+1$, $+5$,
$+15$, and $+30$ minutes. The seven focal tests are adjusted by Holm's method
within each threshold definition. A significant negative-horizon coefficient
would indicate anticipation, price-chasing, shared information, or residual
timing confounding rather than an effect of the later flow.

For $h_2>h_1$, subsequent movement is

\begin{equation}
  \text{Reversal}_{m,t}(h_1,h_2)=
  \operatorname{sign}(\text{NetSignedFlow}_{m,t})
  \left(p_{m,t+h_2}-p_{m,t+h_1}\right).
\end{equation}

A negative value denotes movement against the initial flow direction. A positive value is termed subsequent same-direction continuation, not strict persistence of the initial price movement: it can reflect delayed adjustment when no initial update occurred. A separate descriptive specification therefore conditions on a non-zero initial movement from minute zero to minute five. This conditional sample is outcome-selected and is not used as the primary H3a test. Reversal is not, by itself, proof of manipulation. Let $y_m\in\{0,1\}$ be the resolved outcome. Forecast improvement is measured by the change in Brier loss,

\begin{equation}
  \text{BrierImprovement}_{m,t}(h)=
  (y_m-p_{m,t})^2-(y_m-p_{m,t+h})^2.
\end{equation}

A positive value means that the later probability is closer to the realised outcome. This avoids an inappropriate unconditional 50 per cent accuracy benchmark. Movement to final zero-or-one settlement is not labelled permanent impact because settlement is mechanically determined by resolution.

\subsection{Empirical specifications}

The primary specification regresses the strictly aligned five-minute
pre-kickoff price change on signed-log flow. H1 uses total net signed flow. H2
decomposes flow into past-only P99 large-trade and ordinary-trade components and
directly tests $H_0:\beta_{large}=\beta_{ordinary}$. Primary controls are the
pre-period probability, lagged 30-minute price change, log lagged 30-minute
volume, and log minutes to kickoff. Models include market and unique UTC
calendar-hour fixed effects. Standard errors are clustered by the 104 match
events.

Two additional specifications address timing and shared match information.
First, the pre-trade placebo and post-trade lead--lag models estimate the same
coefficient contrast at all seven frozen horizons. Second, market fixed effects
are combined with event-by-minute and event-by-five-minute fixed effects, so
identification comes from related home-win, draw, and away-win contracts for the
same match in the same time cell. Cells containing only one observed market do
not identify this comparison and are excluded.

These event-by-time fixed effects absorb information common to related
contracts within the same match-time cell, thereby reducing common match-time
confounding. They cannot eliminate contract-specific responses to
contemporaneous match news. The resulting estimates therefore remain
associational.

The extensive margin is estimated first as a linear-probability model for any
non-zero five-minute update and then as a market-fixed-effects logit robustness
model. Markets without within-market outcome variation carry no finite logit
slope information and are excluded from that robustness sample. The intensive
margin conditions on a non-zero update and is treated as descriptive because it
selects on the realised outcome.

H3a is estimated on movement from minute 5 to minute 15; H3b is estimated on the associated Brier-score improvement. RQ4 uses five mutually exclusive pre-kickoff timing bands and a pooled five-minute pre-/post-kickoff phase-interaction model. The phase comparison uses controls common to both phases and tests the change in the large-minus-ordinary coefficient gap. Because RQ4 is exploratory, no monotonic pattern toward kickoff is imposed or expected a priori.

Primary standard errors are clustered by match event because the home-win, away-win, and draw contracts for the same FIFA match are related. Robustness considers market clustering, two-way event and calendar-hour clustering where supported, and non-overlapping five-minute observations. The clustering rule and number of clusters are stated for every table.

\subsection{Secondary wallet-heterogeneity design}
\label{sec:wallet-heterogeneity-design}

Secondary wallet tests distinguish an unusually large individual trade from
past observed-wallet activity and cumulative gross trading volume. High activity
is classified daily using a rolling seven-day top one per cent with minimum
history requirements; expanding history is a robustness definition. Four
large/ordinary by high/other activity cells are compared using price and
sequence outcomes. Complete definitions, HA1--HA4, tests, and multiplicity rules
are reported in Appendix~\ref{sec:wallet-methods-appendix}.

\subsection{Past-only whale-wallet design}
\label{sec:whale-wallet-design}

The past-only whale proxy instead ranks wallets by cumulative gross trading volume before
each evaluation day and imposes no minimum trade-frequency requirement. Rolling
windows, alternative percentiles, and continuous standardised cumulative volume
provide robustness. These measures describe prior observed gross trading volume,
not ownership, wealth, institutional status, or intent. Detailed WH1--WH4 and
outcome definitions are also provided in Appendix~\ref{sec:wallet-methods-appendix}.

\subsection{Identification, robustness, and data limitations}

Large trades may respond to public news or match events, creating simultaneity
and selection. Fixed effects, strict temporal ordering, pre-trade placebos, and
event-by-time comparisons reduce specific forms of confounding but do not
establish causality. The empirical language therefore remains associative.

\subsection{Reproducibility}

The public repository documents the final analysis workflow and contains a
README, stage-specific dependency files, the data dictionary and collection
report, analysis scripts, the principal non-confidential v5 and v6 output
files, SHA-256 manifests for 360 trade partitions and 360 price partitions,
and the complete LaTeX source. Large source and regression-ready datasets are
not redistributed in the repository. Their hash manifests identify the local
files used for the analysis after the public-source data have been collected.

The repository version corresponding to this submitted dissertation is
available at:
\url{https://github.com/andrew1004aa/polymarket-world-cup-2026/tree/v1.0.1-submission}.
Files retained for transparency or development history are not, merely by
their presence, part of the final research design described in this
dissertation.

Detailed specifications for the supplementary concentration,
reversal-prediction, wallet-sequence, and anomaly-screening analyses are reported in
Appendix~\ref{sec:supplementary-integrity-methods}. Keeping those specifications
outside the main chapter preserves the distinction between the primary H1--H3
design and exploratory behavioural or integrity analyses.
.
% Verified references use natbib; remaining legacy citations await metadata verification.

\section{Research Design}
\label{sec:research-design}

\subsection{Research objective and setting}

This study examines whether unusually large individual trades are associated with short-horizon movements in Polymarket-implied probabilities, whether additional flow-aligned movement follows between minutes 5 and 15 or instead reverses, and whether prices improve alignment with eventual market outcomes. The design is observational: it estimates conditional associations and does not claim that a trade caused a price change, conveyed private information, or manipulated a market. The primary outcome is described as a recorded-price adjustment, and the estimated coefficients as directional price associations.

The observation window spans 2026-06-01T00:00:00Z to 2026-08-01T00:00:00Z, with the end point excluded. The final universe consists of 360 binary markets: 312 home-win, away-win, or draw markets covering 104 FIFA World Cup matches, and 48 outright country-winner markets. The canonical trade sample contains 6,725,732 records from the public Polymarket Data API and 335,414 normalised proxy-wallet addresses. The price dataset contains 17,980,605 raw CLOB price-history points for 720 YES and NO tokens, requested at one-minute fidelity without interpolation or resampling. Match timing and outcomes come from the reviewed 104-event FIFA table. Independently obtained Dune Analytics market-level counts were used only to validate extraction completeness; Dune observations were not used to add, remove, or replace canonical trade observations.

\subsection{Research questions and hypotheses}

The study asks four questions. First, as a baseline check, how is signed trade flow associated with subsequent changes in a market's YES-equivalent probability? Second, and central to the dissertation, does unusually large net flow have a stronger association than ordinary net flow? Update incidence and conditional magnitude are examined as mechanism diagnostics rather than embedded in H2. Third, is large-trade flow associated with additional flow-aligned movement between minutes 5 and 15 or reversal over that interval, and do prices improve relative to the realised binary outcome? Fourth, does the association vary with time to kickoff and match phase?

H1 and H2 are primary. Within that pair, H1 is a baseline validation of the signed-flow construction, while H2, together with its extensive-versus-intensive-margin decomposition, is the dissertation's main empirical contribution. H3a, H3b and RQ4 are secondary hypotheses and a research question that interpret and add heterogeneity context to the H2 result rather than stand as co-equal contributions. The corresponding hypotheses and research question are:

\begin{description}
  \item[H1 (baseline).] Greater signed total trade flow is associated with a short-horizon probability change in the same direction. This result validates the signed-flow construction; it is not, by itself, the dissertation's central finding.
  \item[H2 (primary contribution).] Conditional on the common controls and fixed effects, signed flow from individual trades classified by a past-only, market-specific P99 threshold has a larger five-minute directional-price association than signed flow from ordinary trades. Whether the difference operates through update incidence or conditional update magnitude is evaluated as a mechanism decomposition rather than built into the hypothesis.
  \item[H3a.] Large-trade flow is associated with subsequent same-direction continuation from minute five to minute fifteen rather than reversal over that later interval.
  \item[H3b.] Large-trade-associated price movements improve forecast accuracy relative to the eventual resolved outcome, measured by Brier-score improvement, relative to ordinary-trade-associated movements.
\end{description}

The literature reviewed in Section~\ref{sec:literature-review} does not yield a strong ex-ante directional prediction for the shape of pre-kickoff timing heterogeneity, so this dimension is treated as a research question rather than forced into a falsifiable directional hypothesis:

\begin{description}
  \item[RQ4.] How does the large-versus-ordinary flow--price association vary across pre-specified time-to-kickoff bands, and between the pre-kickoff and post-kickoff/pre-resolution phases?
\end{description}

This formulation follows the microstructure literature on informed trading, trade size, and price discovery \citep{kyle1985continuous,easleyohara1987tradesize,hasbrouck1991information,madhavan1997prices}, while prediction-market research motivates the interpretation of prices as market-implied probabilities subject to institutional and belief-based qualifications \citep{wolferszitzewitz2004prediction,manski2006interpreting}. Manipulation and corrective-trading research motivates the distinction between temporary pressure and information \citep{allengale1992manipulation,camerer1998racetrack,hanson2006aggregation}.

\subsection{Observation units and sample separation}

The primary observation unit is the market-minute. For a minute beginning at
$t$, flow contains trades timestamped in $[t,t+1)$. The baseline is the latest
observed YES-equivalent price strictly before $t$, and each post-flow outcome is
the first observed price at or after the designated horizon. This ordering
prevents the baseline from incorporating a price observation recorded after a
trade in the focal minute.

The match and outright families are estimated separately. The primary match specification uses pre-kickoff observations. Post-kickoff/pre-resolution observations are analysed separately because match events generate fundamental information; this phase is not labelled purely in-play because actual final-whistle timestamps are unavailable. Outright markets form a separate panel because their horizon and information environment differ from single-match contracts. Trades at or after recorded resolution remain in the method-neutral audit panel but are excluded from model samples.

\subsection{Direction and large-trade definition}

Trade direction is expressed in YES-probability terms. BUY YES and SELL NO receive a sign of $+1$; SELL YES and BUY NO receive a sign of $-1$. For trade $i$,

\begin{equation}
  \text{SignedValue}_{i}=\text{Direction}_{i}\times
  \left(\text{ExecutedTokenPrice}_{i}\times\text{SizeShares}_{i}\right).
\end{equation}

A YES-token history is retained as observed, whereas a NO-token price is converted to one minus the NO price. YES and NO histories are not treated as independent markets. The YES history is primary and the NO-equivalent series is used for complementarity and data-quality checks.

The CLOB history documentation labels each observation only as a timestamp and
price value; it does not state that the value is a contemporaneous midpoint or
pre-trade quote. The primary dependent variable is therefore termed a
recorded-price adjustment. A robustness series instead uses the last executed
YES-equivalent trade price strictly before the focal minute and the first
executed trade price at or after the target horizon, with 60-, 300-, and
900-second matching tolerances. The focal minute is excluded from the
executed-price baseline.

For each binary market, trade value is executed token price multiplied by
shares. The primary expanding-history threshold for UTC day $d$ is the linearly
interpolated P99 of unsigned USDC trade value in the same market from 1 June to
the end of day $d-1$. A threshold becomes eligible after at least 1,000 prior
trades. The principal definition robustness instead uses the preceding seven
complete UTC days. A trade is large when its value is at or above the applicable
threshold; ties are included. Both definitions are past-only, so a trade is
never classified using later observations. Models are estimated separately and
on their common support. Past-only P95, P97.5, and P99.5 thresholds test whether
the finding depends on the P99 boundary. Signed-log flow is also replaced by
\texttt{asinh} flow and winsorised raw flow scaled in \$1,000 units;
standardised and common-value effects permit comparison despite the different
sparsity and scale of the large and ordinary components. Classification occurs before applying YES-equivalent
direction and does not represent wallet wealth, ownership, institutional status,
profitability, information, or intent.

\subsection{Price response, subsequent continuation, and forecast improvement}

Let $p^-_{m,t}$ be the latest available YES-equivalent CLOB price strictly
before $t$, and let $p^+_{m,t+h}$ be the first observed price at or after the
target $t+h$. Prices are not interpolated. For horizon $h$,

\begin{equation}
  \Delta p_{m,t}(h)=p^+_{m,t+h}-p^-_{m,t}.
\end{equation}

The primary pre-kickoff match horizon is five minutes; one, 15, and 30 minutes
trace the response path. Strict-0 alignment requires the baseline to occur less
than 60 seconds before $t$ and the outcome less than 60 seconds after its target;
strict-1 permits 120 seconds. No interpolation, resampling, or gap filling is
used. No pre-kickoff model window may cross actual kickoff.

Pre-trade placebo outcomes cover $[-15,0)$, $[-5,0)$, and $[-1,0)$ relative
to the focal flow minute. Post-trade lead--lag outcomes cover $+1$, $+5$,
$+15$, and $+30$ minutes. The seven focal tests are adjusted by Holm's method
within each threshold definition. A significant negative-horizon coefficient
would indicate anticipation, price-chasing, shared information, or residual
timing confounding rather than an effect of the later flow.

For $h_2>h_1$, subsequent movement is

\begin{equation}
  \text{Reversal}_{m,t}(h_1,h_2)=
  \operatorname{sign}(\text{NetSignedFlow}_{m,t})
  \left(p_{m,t+h_2}-p_{m,t+h_1}\right).
\end{equation}

A negative value denotes movement against the initial flow direction. A positive value is termed subsequent same-direction continuation, not strict persistence of the initial price movement: it can reflect delayed adjustment when no initial update occurred. A separate descriptive specification therefore conditions on a non-zero initial movement from minute zero to minute five. This conditional sample is outcome-selected and is not used as the primary H3a test. Reversal is not, by itself, proof of manipulation. Let $y_m\in\{0,1\}$ be the resolved outcome. Forecast improvement is measured by the change in Brier loss,

\begin{equation}
  \text{BrierImprovement}_{m,t}(h)=
  (y_m-p_{m,t})^2-(y_m-p_{m,t+h})^2.
\end{equation}

A positive value means that the later probability is closer to the realised outcome. This avoids an inappropriate unconditional 50 per cent accuracy benchmark. Movement to final zero-or-one settlement is not labelled permanent impact because settlement is mechanically determined by resolution.

\subsection{Empirical specifications}

The primary specification regresses the strictly aligned five-minute
pre-kickoff price change on signed-log flow. H1 uses total net signed flow. H2
decomposes flow into past-only P99 large-trade and ordinary-trade components and
directly tests $H_0:\beta_{large}=\beta_{ordinary}$. Primary controls are the
pre-period probability, lagged 30-minute price change, log lagged 30-minute
volume, and log minutes to kickoff. Models include market and unique UTC
calendar-hour fixed effects. Standard errors are clustered by the 104 match
events.

Two additional specifications address timing and shared match information.
First, the pre-trade placebo and post-trade lead--lag models estimate the same
coefficient contrast at all seven frozen horizons. Second, market fixed effects
are combined with event-by-minute and event-by-five-minute fixed effects, so
identification comes from related home-win, draw, and away-win contracts for the
same match in the same time cell. Cells containing only one observed market do
not identify this comparison and are excluded.

These event-by-time fixed effects absorb information common to related
contracts within the same match-time cell, thereby reducing common match-time
confounding. They cannot eliminate contract-specific responses to
contemporaneous match news. The resulting estimates therefore remain
associational.

The extensive margin is estimated first as a linear-probability model for any
non-zero five-minute update and then as a market-fixed-effects logit robustness
model. Markets without within-market outcome variation carry no finite logit
slope information and are excluded from that robustness sample. The intensive
margin conditions on a non-zero update and is treated as descriptive because it
selects on the realised outcome.

H3a is estimated on movement from minute 5 to minute 15; H3b is estimated on the associated Brier-score improvement. RQ4 uses five mutually exclusive pre-kickoff timing bands and a pooled five-minute pre-/post-kickoff phase-interaction model. The phase comparison uses controls common to both phases and tests the change in the large-minus-ordinary coefficient gap. Because RQ4 is exploratory, no monotonic pattern toward kickoff is imposed or expected a priori.

Primary standard errors are clustered by match event because the home-win, away-win, and draw contracts for the same FIFA match are related. Robustness considers market clustering, two-way event and calendar-hour clustering where supported, and non-overlapping five-minute observations. The clustering rule and number of clusters are stated for every table.

\subsection{Secondary wallet-heterogeneity design}
\label{sec:wallet-heterogeneity-design}

Secondary wallet tests distinguish an unusually large individual trade from
past observed-wallet activity and cumulative gross trading volume. High activity
is classified daily using a rolling seven-day top one per cent with minimum
history requirements; expanding history is a robustness definition. Four
large/ordinary by high/other activity cells are compared using price and
sequence outcomes. Complete definitions, HA1--HA4, tests, and multiplicity rules
are reported in Appendix~\ref{sec:wallet-methods-appendix}.

\subsection{Past-only whale-wallet design}
\label{sec:whale-wallet-design}

The past-only whale proxy instead ranks wallets by cumulative gross trading volume before
each evaluation day and imposes no minimum trade-frequency requirement. Rolling
windows, alternative percentiles, and continuous standardised cumulative volume
provide robustness. These measures describe prior observed gross trading volume,
not ownership, wealth, institutional status, or intent. Detailed WH1--WH4 and
outcome definitions are also provided in Appendix~\ref{sec:wallet-methods-appendix}.

\subsection{Identification, robustness, and data limitations}

Large trades may respond to public news or match events, creating simultaneity
and selection. Fixed effects, strict temporal ordering, pre-trade placebos, and
event-by-time comparisons reduce specific forms of confounding but do not
establish causality. The empirical language therefore remains associative.

\subsection{Reproducibility}

The public repository documents the final analysis workflow and contains a
README, stage-specific dependency files, the data dictionary and collection
report, analysis scripts, the principal non-confidential v5 and v6 output
files, SHA-256 manifests for 360 trade partitions and 360 price partitions,
and the complete LaTeX source. Large source and regression-ready datasets are
not redistributed in the repository. Their hash manifests identify the local
files used for the analysis after the public-source data have been collected.

The repository version corresponding to this submitted dissertation is
available at:
\url{https://github.com/andrew1004aa/polymarket-world-cup-2026/tree/v1.0.1-submission}.
Files retained for transparency or development history are not, merely by
their presence, part of the final research design described in this
dissertation.

Detailed specifications for the supplementary concentration,
reversal-prediction, wallet-sequence, and anomaly-screening analyses are reported in
Appendix~\ref{sec:supplementary-integrity-methods}. Keeping those specifications
outside the main chapter preserves the distinction between the primary H1--H3
design and exploratory behavioural or integrity analyses.
